\chapter{\IfLanguageName{dutch}{Stand van zaken}{State of the art}}%
\label{ch:stand-van-zaken}

Dit hoofdstuk bouwt verder op de context en probleemstelling uit Hoofdstuk~\ref{ch:inleiding} door de relevante literatuur rond IP-camera’s, IoT-beveiliging en netwerkintegratie te bespreken. Het doel is om de belangrijkste risico’s, aanvalsvectoren en gangbare mitigaties te situeren, en zo een onderbouwd kader te creëren voor de methodologische keuzes in Hoofdstuk~\ref{ch:methodologie}. De literatuurstudie focust achtereenvolgens op (1) de plaats van IP-camera’s binnen het IoT-domein, (2) typische kwetsbaarheden, (3) de impact van netwerkintegratie op het aanvalsoppervlak, (4) contextfactoren in retail- en KMO-omgevingen, (5) geschikte beveiligingskaders en evaluatiemodellen, en (6) de onderzoekslacune die de aanleiding vormt voor deze bachelorproef.

\section{\IfLanguageName{dutch}{IP-camera’s binnen het IoT-domein}{IP cameras in the IoT domain}}%
\label{sec:camera-iot}

IP-camera’s maken deel uit van het bredere Internet of Things (IoT)-ecosysteem, waarin fysieke apparaten via netwerken gegevens verzamelen, verwerken en uitwisselen. IoT-systemen worden doorgaans gemodelleerd volgens een gelaagde architectuur, waarbij beveiligingsuitdagingen zich manifesteren op meerdere niveaus, gaande van fysieke apparaten tot applicatie- en businesslagen \autocite{Butun2019arXivIoTSecurity,rajmohandecade}. Binnen dergelijke architecturen vormen IoT-apparaten op de perceptie- of device-laag vaak de zwakste schakel.

Onderzoek toont aan dat IoT-systemen een aanzienlijk vergroot aanvalsoppervlak introduceren door hun grootschalige connectiviteit, heterogene protocollen en permanente netwerktoegang \autocite{MOHAMADNOOR2019283,rajmohandecade}. IP-camera’s functioneren als continu verbonden embedded systemen binnen bedrijfsnetwerken en cloudomgevingen, waardoor zij potentiële toegangspunten vormen voor laterale bewegingen binnen een infrastructuur.

Een fundamenteel kenmerk van IoT-apparaten is hun resourcebeperkte aard. In tegenstelling tot traditionele IT-systemen beschikken zij over beperkte rekenkracht, geheugen en energiecapaciteit, wat de implementatie van conventionele beveiligingsmechanismen bemoeilijkt \autocite{fysarakis2014embedded,kumar2024review}. Hierdoor ontstaat een spanningsveld tussen performantie, kostenefficiëntie en beveiliging. Zwaardere cryptografische protocollen of uitgebreide loggingmechanismen kunnen de functionaliteit of responstijd van het apparaat negatief beïnvloeden.

Daarnaast blijkt uit literatuur dat IoT-beveiligingsonderzoek sterk focust op architecturale oplossingen, terwijl specifieke beveiligingspatronen op device-niveau ondervertegenwoordigd blijven \autocite{rajmohandecade}. Deze discrepantie suggereert dat hoewel high-level beveiligingsarchitecturen worden voorgesteld, de praktische implementatie op het niveau van individuele apparaten, zoals IP-camera’s, vaak onvoldoende uitgewerkt blijft.

Om deze beperkingen te mitigeren, wordt in recente literatuur gewezen op het belang van lightweight cryptografische oplossingen en contextspecifieke beveiligingsmechanismen die aangepast zijn aan resource-constrained omgevingen \autocite{kumar2024review}. Dergelijke oplossingen trachten een evenwicht te vinden tussen cryptografische sterkte en operationele haalbaarheid, zonder de prestaties van het IoT-apparaat significant te degraderen.

Samenvattend situeren IP-camera’s zich binnen een complex IoT-ecosysteem waarin architecturale, technische en resourcegerelateerde beperkingen samenkomen. De combinatie van permanente netwerkconnectiviteit, beperkte hardwarecapaciteit en vaak gebrekkige secure-by-design implementatie maakt deze apparaten bijzonder relevant binnen hedendaags IoT-beveiligingsonderzoek \autocite{MOHAMADNOOR2019283,rajmohandecade}.

\section{\IfLanguageName{dutch}{Typische kwetsbaarheden bij IP-camera’s}{Typical vulnerabilities in IP cameras}}%
\label{sec:typische-kwetsbaarheden}

Hoewel IP-camera’s primair ontworpen zijn voor videobewaking en monitoring, vertonen zij in de praktijk terugkerende kwetsbaarheden die systematisch worden gedocumenteerd binnen het IoT- en embedded security-onderzoek. Deze kwetsbaarheden situeren zich voornamelijk op het niveau van authenticatie, firmwarebeheer, netwerkconfiguratie en communicatieprotocollen \autocite{baho2023analysis,Butun2019arXivIoTSecurity,ENISA2024ThreatLandscape}. Binnen gelaagde IoT-architecturen bevinden veel van deze zwaktes zich op de fysieke of device-laag, waar beperkte beveiligingsmechanismen gecombineerd worden met hoge netwerkexpositie \autocite{baho2023analysis}.

\subsection{\IfLanguageName{dutch}{Authenticatie en standaardinloggegevens}{Authentication and default credentials}}%
\label{subsec:auth-defaults}

Een van de meest voorkomende problemen betreft het gebruik van zwakke of onveranderde standaardinloggegevens. Veel IP-camera’s worden geleverd met vooraf ingestelde gebruikersnamen en wachtwoorden die in productieomgevingen niet systematisch worden gewijzigd. Zowel academische literatuur als praktijkrapporten wijzen erop dat hardcoded of eenvoudig te raden credentials een structureel risico vormen binnen IoT-apparaten \autocite{baho2023analysis,Butun2019arXivIoTSecurity}. De grootschalige exploitatie van dergelijke zwaktes werd onder meer zichtbaar in botnetcampagnes zoals Mirai, waarbij DVR- en IP-camerasystemen werden gecompromitteerd via bekende standaardcredentials \autocite{zhou2023survey}. Dit illustreert dat het probleem niet louter theoretisch is, maar reeds op grote schaal operationeel werd misbruikt.

\subsection{\IfLanguageName{dutch}{Webgebaseerde beheerinterfaces}{Web-based management interfaces}}%
\label{subsec:web-interfaces}

IP-camera’s bieden doorgaans configuratietoegang via webinterfaces die draaien op embedded webservers. Onderzoek naar embedded device security toont aan dat dergelijke interfaces vaak kwetsbaar zijn voor brute-force-aanvallen, command injection, cross-site scripting of onvoldoende sessiebeheer \autocite{zhou2023survey}. Door de afwezigheid van geavanceerde beveiligingsmechanismen zoals Address Space Layout Randomization (ASLR) of Control-Flow Integrity in veel embedded systemen, zijn deze webinterfaces bijzonder gevoelig voor exploitatie \autocite{zhou2023survey}. Wanneer dergelijke interfaces rechtstreeks toegankelijk zijn vanuit bredere netwerksegmenten, vergroot dit aanzienlijk het aanvalsoppervlak van de organisatie.

\subsection{\IfLanguageName{dutch}{Onversleutelde communicatie en protocolrisico’s}{Unencrypted communication and protocol risks}}%
\label{subsec:protocol-risks}

IP-camera’s maken vaak gebruik van protocollen zoals RTSP of HTTP voor het verzenden van videostreams. Indien deze communicatie niet adequaat versleuteld wordt, kunnen netwerkgebaseerde aanvallen zoals sniffing of man-in-the-middle-aanvallen leiden tot onderschepping of manipulatie van videodata \autocite{ENISA2024ThreatLandscape}. Kwetsbaarheden in communicatie- en netwerkprotocollen worden bovendien frequent geïdentificeerd binnen IoT-ecosystemen, waarbij implementatiefouten of configuratiegebreken bijdragen aan verhoogde blootstelling \autocite{baho2023analysis}. Het ontbreken van netwerksegmentatie kan dit risico verder versterken.

\subsection{\IfLanguageName{dutch}{Firmwarekwetsbaarheden en update-mechanismen}{Firmware vulnerabilities and update mechanisms}}%
\label{subsec:firmware-updates}

Firmwarebeheer vormt een kritieke component binnen de beveiliging van IP-camera’s. Embedded devices bevatten vaak meerdere kritieke kwetsbaarheden in firmwarecomponenten, terwijl patchbeheer en updatecycli in de praktijk onregelmatig of onvoldoende gecontroleerd verlopen \autocite{zhou2023survey,ENISA2024ThreatLandscape}. Recent onderzoek naar firmware-updateprocessen toont bovendien aan dat update-mechanismen zelf kwetsbaar kunnen zijn. Ibrahim et al.\ identificeren firmware modification attacks, downgrade attacks en device bricking attacks als gevolg van ontbrekende cryptografische verificatie of rollback-preventie binnen IoT-updateprocessen \autocite{ibrahim2023aot}. Dergelijke kwetsbaarheden impliceren dat niet alleen verouderde firmware een risico vormt, maar ook het updateproces zelf een aanvalsvector kan zijn. Het bestaan van publiek gedocumenteerde kwetsbaarheden in de National Vulnerability Database (NVD) wijst erop dat bekende technische zwaktes langdurig aanwezig kunnen blijven in operationele systemen wanneer updates niet tijdig worden toegepast \autocite{NVD2025}.

\subsection{\IfLanguageName{dutch}{Interoperabiliteit en configuratiecomplexiteit}{Interoperability and configuration complexity}}%
\label{subsec:interoperabiliteit}

Interoperabiliteitstandaarden zoals ONVIF bevorderen compatibiliteit tussen cameramodellen en beheersystemen \autocite{ONVIFCoreSpec2025}. Hoewel standaardisatie operationele voordelen biedt, kan foutieve implementatie of configuratie van toegangsrechten leiden tot onbedoelde blootstelling van functionaliteiten. Onderzoek naar IoT-kwetsbaarheden benadrukt bovendien dat beveiligingsproblemen zelden geïsoleerd optreden, maar vaak voortkomen uit een combinatie van ontwerplimieten, configuratiefouten en organisatorische tekortkomingen \autocite{baho2023analysis,Butun2019arXivIoTSecurity}.

\subsection{\IfLanguageName{dutch}{Synthese}{Synthesis}}%
\label{subsec:synthese-kwetsbaarheden}

De literatuur toont aan dat kwetsbaarheden bij IP-camera’s zich niet beperken tot individuele technische defecten, maar voortvloeien uit de interactie tussen resourcebeperkingen, embedded architectuur, updatebeheer en netwerkconfiguratie \autocite{Butun2019arXivIoTSecurity,zhou2023survey}. Hoewel bestaande studies een breed overzicht bieden van typische aanvalsvectoren binnen IoT-omgevingen, blijven contextspecifieke analyses van heterogene IP-camerainfrastructuren in bedrijfsomgevingen relatief beperkt. Dit suggereert dat een empirische, contextgebonden evaluatie noodzakelijk is om de praktische impact van deze kwetsbaarheden adequaat te beoordelen.

\section{\IfLanguageName{dutch}{Netwerkintegratie en attack surface uitbreiding}{Network integration and attack surface expansion}}%
\label{sec:netwerkintegratie-attack-surface}

De beveiligingsrisico’s van IP-camera’s kunnen niet los worden gezien van hun integratie binnen bredere netwerkinfrastructuren. Hoewel kwetsbaarheden op toestel- of firmware-niveau relevant blijven, wordt hun impact in sterke mate bepaald door de netwerkarchitectuur waarin camera’s opereren. In IoT-omgevingen wordt herhaaldelijk benadrukt dat elk bijkomend verbonden apparaat het totale aanvalsoppervlak vergroot en extra interfaces en beheercomplexiteit introduceert \autocite{Butun2019arXivIoTSecurity}. Een systematische manier om dat aanvalsoppervlak te benaderen is het identificeren van alle relevante componenten, interfaces en datastromen over de volledige levenscyclus van een toestel, omdat net daar vaak onverwachte aanvalspaden ontstaan \autocite{Liebl2024arXivIIoTAttackSurface}.

IP-camera’s worden in de praktijk vaak geïntegreerd in interne LAN-omgevingen waar ze communiceren met Network Video Recorders (NVR’s), videomanagementsoftware en soms cloud- of remote access-diensten. Wanneer netwerksegmentatie ontbreekt of beperkt wordt toegepast, kan een kwetsbaarheid in één camera leiden tot laterale beweging binnen het netwerk. ENISA wijst netwerksegmentatie expliciet aan als terugkerend pijnpunt wanneer IoT-apparaten worden toegevoegd zonder expliciete beveiligingsarchitectuur \autocite{ENISA2024ThreatLandscape}. Bovendien tonen recente enterprise-studies rond lateral movement dat deze fase kernachtig is binnen APT-aanvallen en dat betrouwbare detectie in complexe, evoluerende netwerken moeilijk blijft, waardoor preventieve beperking (bv.\ segmentatie en least privilege) extra belangrijk wordt \autocite{khoury2024jbeil}.

Beveiligingskaders zoals NIST SP 800-53 benadrukken daarom het belang van netwerkisolatie, toegangscontrole en monitoring om de impact van compromittering te beperken \autocite{NISTSP80053r5}. In de praktijk sluit dit aan bij Zero Trust-principes waarbij microsegmentatie wordt ingezet om de \emph{blast radius} te verkleinen en laterale beweging te bemoeilijken \autocite{wonor2025zero}.

Een bijkomende complexiteit ontstaat wanneer meerdere cameratypes of -merken binnen één infrastructuur worden gecombineerd. Interoperabiliteit via standaarden zoals ONVIF maakt integratie tussen verschillende toestellen mogelijk \autocite{ONVIFCoreSpec2025}, maar verhoogt tegelijkertijd de complexiteit van toegangsbeheer, roltoewijzing en configuratieconsistentie. Heterogene omgevingen leiden vaker tot uiteenlopende firmwareversies, verschillende default-instellingen en variërende implementaties van authenticatie- en updateprocessen. Net die variatie vergroot de kans op misconfiguraties, wat in IoT-literatuur consequent als belangrijke risicofactor wordt aangeduid \autocite{Butun2019arXivIoTSecurity}.

Daarnaast wordt in de praktijk vaak externe toegang tot camerabeelden voorzien, bijvoorbeeld via port forwarding, VPN-verbindingen of cloudgebaseerde diensten. Hoewel dergelijke functionaliteiten operationeel nuttig zijn, introduceren ze bijkomende blootstelling aan het internet. Dit is relevant omdat internet-wide scanning een wijdverspreide en schaalbare techniek is waarmee kwetsbare diensten zeer snel opgespoord kunnen worden, zeker wanneer nieuwe kwetsbaarheden publiek worden \autocite{durumeric2014internet}. ENISA benadrukt dat foutief geconfigureerde externe toegangsmogelijkheden regelmatig aanleiding geven tot misbruik van IoT-apparaten \autocite{ENISA2024ThreatLandscape}. Wanneer basismaatregelen (bv.\ sterke authenticatie, beperking van bron-IP’s, minimale blootstelling van services) ontbreken, stijgt de kans op ongeautoriseerde toegang en verdere pivoting aanzienlijk.

De literatuur suggereert bijgevolg dat de beveiligingspositie van IP-camera’s niet uitsluitend bepaald wordt door de intrinsieke eigenschappen van het toestel, maar door de interactie tussen toestel, configuratie en netwerkarchitectuur \autocite{Liebl2024arXivIIoTAttackSurface,NISTSP80053r5}. Dit onderstreept het belang van een contextgebonden analyse waarin niet alleen toestelkwetsbaarheden, maar ook netwerkstructuur, toegangscontrole en externe exposure systematisch worden geëvalueerd.

\section{\IfLanguageName{dutch}{Specifieke risico’s binnen retail- en KMO-omgevingen}{Specific risks in retail and SME environments}}%
\label{sec:retail-kmo}

Hoewel de literatuur uitgebreid ingaat op IoT-beveiliging en kwetsbaarheden van netwerkapparatuur, krijgt de organisatorische context waarin deze systemen worden beheerd vaak minder aandacht. In retail- en KMO-omgevingen verschilt de implementatie en het beheer van IT-infrastructuur doorgaans sterk van grotere enterprise-omgevingen, waardoor dezelfde technische kwetsbaarheid in de praktijk een andere impact kan hebben. ENISA benadrukt dat IoT- en netwerkapparatuur terugkerende doelwitten zijn binnen het actuele dreigingslandschap, maar dat de weerbaarheid sterk wordt beïnvloed door factoren zoals governance, patchdiscipline en basis-hygiëne in configuratie \autocite{ENISA2024ThreatLandscape}.

Binnen kleine en middelgrote ondernemingen is IT-beheer vaak een ondersteunende taak, waarbij in retailomgevingen operationele continuïteit primeert boven formele IT-architectuur. ENISA wijst hierbij op beperkte middelen, gebrek aan gespecialiseerde beveiligingsexpertise en het ontbreken van formele security governance als structurele risicofactoren \autocite{ENISA2024ThreatLandscape}. In de context van IP-camera’s vertaalt dit zich naar een verhoogde kans dat toestellen wel functioneel worden geïntegreerd, maar dat structurele processen zoals patchbeheer, periodieke configuratiereviews en monitoring niet consequent worden toegepast.

Een bijkomende uitdaging betreft de organische groei van infrastructuur. Wanneer uitbreidingen over langere tijd plaatsvinden, ontstaat vaak een heterogene omgeving met verschillende toestellen, firmwareversies en beheerinterfaces. Dit verhoogt de beheercomplexiteit en daarmee de kans op inconsistente beveiligingsinstellingen en misconfiguraties \autocite{Butun2019arXivIoTSecurity}. Empirische analyses tonen bovendien dat in grote, heterogene netwerken IoT-apparaten (waaronder CCTV/IP-camera’s) geregeld draaien met verouderde firmware of default credentials, mede door gebrekkige zichtbaarheid op device-inventaris en eigenaarschap \autocite{agarwal2019detecting}.

Naast technische factoren speelt ook risicoperceptie een rol. In kleinere organisaties wordt camerabewaking vaak primair benaderd als maatregel voor fysieke veiligheid en diefstalpreventie, waardoor het digitale aanvalsvlak minder zichtbaar wordt in besluitvorming en beheerprioriteiten \autocite{ENISA2024ThreatLandscape}. Vanuit onderzoek naar risicoperceptie bij IoT blijkt bovendien dat het communiceren van security- en privacy-eigenschappen wel degelijk beïnvloedt hoe risico’s worden ingeschat, maar dat risicoperceptie en feitelijk gedrag niet altijd perfect aligneren \autocite{emami2021privacy}.

Tot slot tonen internet-scale metingen dat misconfiguraties en zwakke instellingen bij IoT-apparaten geen uitzonderingen zijn, maar op grote schaal voorkomen. Srinivasa et al.\ identificeren miljoenen misgeconfigureerde IoT-devices en laten zien dat aanvallers misconfiguraties actief misbruiken in het wild \autocite{srinivasa2021open}. Dit is bijzonder relevant voor retail- en KMO-contexten, waar tijdsdruk, beperkte expertise en heterogene installaties het risico op configuratiefouten verhogen.

Hoewel bestaande studies waardevolle inzichten bieden in IoT-kwetsbaarheden en netwerkbeveiliging, blijven ze vaak gericht op grootschalige infrastructuren of generieke modellen. Retailgerichte literatuur benadrukt bovendien dat cybersecurity-implementatie in de sector vaak moet balanceren tussen beveiligingsnoden en operationele haalbaarheid, wat de adoptie van structurele maatregelen kan vertragen \autocite{ardhaninggar2024review}. Dit suggereert dat er nood is aan praktijkgerichte analyses waarin technische kwetsbaarheden worden geëvalueerd binnen hun specifieke organisatorische en operationele context.

\section{\IfLanguageName{dutch}{Beveiligingskaders en evaluatiemodellen}{Security frameworks and evaluation models}}%
\label{sec:kaders-modellen}

De analyse van beveiligingsrisico’s binnen IP-camerainfrastructuren vereist een gestructureerde benadering die verder gaat dan het identificeren van individuele kwetsbaarheden. Bij IP-camera’s ontstaat risico vaak door de combinatie van toestelkenmerken, configuratie, beheerprocessen en netwerkintegratie, waardoor een methodisch kader nodig is om dreigingen systematisch te herkennen en te vertalen naar haalbare mitigaties. Dit is extra relevant in omgevingen waar basiscontroles niet consequent worden toegepast: in KMO-contexten wordt het nut van minimale, prioritaire control-sets benadrukt om de kloof tussen beschikbare standaarden en praktische implementatie te overbruggen \autocite{pawar2022lcci}.

Een veelgebruikte methode binnen threat modeling is het STRIDE-model, ontwikkeld door Microsoft. STRIDE categoriseert dreigingen in zes categorieën: spoofing, tampering, repudiation, information disclosure, denial of service en elevation of privilege \autocite{Microsoft2022bSTRIDE}. Dat STRIDE ook buiten klassieke softwarecontexten bruikbaar is, blijkt uit toepassingen in IoT-scenario’s waar componenten en datastromen expliciet worden gemodelleerd om bedreigingen en mitigaties te inventariseren \autocite{AlAsif2022arXivSTRIDEAgri}. Daarnaast bestaat er literatuur die STRIDE toepast op cyber-physical systems, wat de bruikbaarheid van het model in omgevingen met fysieke componenten ondersteunt \autocite{khan2017stride}.

Naast het identificeren van dreigingen is een controlekader nodig om te bepalen welke beheersmaatregelen passend zijn. NIST SP 800-53 beschrijft een uitgebreide set security- en privacycontroles voor informatiesystemen \autocite{NISTSP80053r5}. Hoewel deze richtlijnen vaak met grotere organisaties worden geassocieerd, blijven kernprincipes zoals toegangscontrole, configuratiebeheer, logging/monitoring en netwerksegmentatie rechtstreeks toepasbaar op camerainfrastructuren in kleinere omgevingen. Voor IoT-apparaten is het bovendien nuttig om naast generieke controles ook device-gerichte baselines te hanteren. NISTIR 8259A specificeert daarom een core baseline van IoT cybersecurity-capabilities (zoals device identification, veilige configuratie, software update-mechanismen en logging), die als referentie kan dienen om toestellen en hun beheerbaarheid objectiever te beoordelen \autocite{fagan2020iot}.

Op Europees niveau benadrukt de ENISA Threat Landscape eveneens dat IoT- en netwerkapparatuur terugkerende doelwitten zijn en wijst ENISA op mitigaties zoals patchbeheer, sterke authenticatie, logging en netwerkisolatie \autocite{ENISA2024ThreatLandscape}. In methodologische termen ontstaat hiermee een bruikbare combinatie: STRIDE helpt om dreigingen te structureren, terwijl NIST/ENISA de vertaalslag ondersteunt naar concrete beheersmaatregelen en prioriteiten.

Voor evaluatie volstaat het echter niet om kwetsbaarheden enkel te inventariseren; het is ook nodig om hun waarschijnlijkheid en impact in context te wegen. In de literatuur rond IIoT-protocollen wordt bijvoorbeeld aangetoond dat CVSS bruikbaar is om kwetsbaarheden te beschrijven, maar dat risico-inschatting afhankelijk blijft van contextfactoren zoals exposure, configuratie en aanwezige mitigaties \autocite{figueroa2020survey}. Dit sluit aan bij de kern van een contextgebonden evaluatie van IP-camera’s: eenzelfde CVE kan in een strikt gesegmenteerd netwerk met beperkte externe toegang en sterke authenticatie een beduidend lagere praktische impact hebben dan in een vlak netwerk met internet-exposure en zwakke beheerinstellingen.

\section{\IfLanguageName{dutch}{Onderzoekslacune}{Research gap}}%
\label{sec:onderzoekslacune}

Uit de voorgaande literatuur blijkt dat de beveiligingsproblematiek rond IoT-apparaten, en in het bijzonder IP-camera’s, uitgebreid wordt besproken binnen zowel academisch onderzoek als professionele dreigingsrapporten. Studies wijzen op structurele zwaktes rond authenticatie, firmwarebeheer en netwerkconfiguratie \autocite{Butun2019arXivIoTSecurity,ENISA2024ThreatLandscape}. Daarnaast benadrukken beveiligingskaders zoals NIST SP 800-53 en threat modeling met STRIDE het belang van gestructureerde dreigingsanalyse, controlemaatregelen en consistente mitigatiekeuzes \autocite{NISTSP80053r5,Microsoft2022bSTRIDE}.

Hoewel deze literatuur waardevolle inzichten biedt in algemene kwetsbaarheden en mitigatiestrategieën, blijven verschillende aspecten onderbelicht wanneer men deze toepast op concrete operationele omgevingen met heterogene cameratoestellen.

Ten eerste focust een aanzienlijk deel van het bestaande onderzoek op generieke IoT-categorieën of op grootschalige infrastructuren, terwijl de dynamiek van heterogene device-ecosystemen in realistische netwerken moeilijker te beheren is. Empirisch onderzoek toont bijvoorbeeld dat IoT-firmware sterk afhankelijk kan zijn van third-party libraries en dat vendors security-updates in de praktijk soms traag doorvoeren, waardoor zelfs up-to-date firmware kwetsbare componenten kan blijven bevatten \autocite{zhang2021capture}. Dit ondersteunt de stelling dat heterogeniteit en software-supply-chain factoren patchbeheer in reële omgevingen complex maken, maar dergelijke aspecten worden minder systematisch onderzocht voor gemengde IP-camerainfrastructuren waarin verschillende modellen en firmwarelijnen naast elkaar bestaan.

Ten tweede blijft de toepassing van algemene IoT-beveiligingsprincipes binnen kleinschalige retail- en KMO-contexten relatief ondergedocumenteerd. Dreigingsrapporten benoemen wel dat IoT- en netwerkapparatuur terugkerende doelwitten zijn en dat basismaatregelen essentieel blijven, maar ze bieden doorgaans minder diepgaande praktijkanalyses van hoe die risico’s zich manifesteren in organisch gegroeide omgevingen waar IT-beheer geen kernactiviteit is \autocite{ENISA2024ThreatLandscape}. Dit sluit aan bij bevindingen dat IoT-devices in grote, heterogene netwerken vaak moeilijk volledig zichtbaar en beheersbaar zijn, wat de kans op verouderde firmware en onveilige instellingen verhoogt \autocite{agarwal2019detecting}.

Ten derde verwijst de literatuur vaak naar bekende kwetsbaarheden en CVE’s, maar minder naar de effectieve exploiteerbaarheid binnen een specifieke operationele context. Het louter bestaan van een gedocumenteerde kwetsbaarheid impliceert niet noodzakelijk dat deze in een concrete infrastructuur ook praktisch misbruikbaar is: exploitability wordt mede bepaald door exposure, configuratie, netwerkpositie en aanwezige mitigaties \autocite{figueroa2020survey,baho2023analysis}. Tegelijk tonen internet-scale metingen dat misconfiguraties en zwakke instellingen bij IoT-devices in de praktijk wél grootschalig voorkomen en actief misbruikt worden, wat de relevantie van gecontroleerde validatie in realistische configuraties versterkt \autocite{srinivasa2021open}.

Ten slotte worden beveiligingskaders zoals STRIDE en NIST SP 800-53 geregeld conceptueel besproken, maar minder vaak toegepast op een concrete, heterogene IP-camerainfrastructuur met als expliciet doel contextgebonden mitigaties te formuleren en te prioriteren. Er bestaan wel voorbeelden waarin STRIDE op componentniveau wordt toegepast om threats te identificeren en te koppelen aan mitigaties \autocite{khan2017stride}, maar de vertaling van die aanpak naar een operationele retailomgeving met heterogene camera’s en pragmatische beperkingen blijft beperkt uitgewerkt. Bovendien vereist empirische validatie via proof-of-concepttests een reproduceerbare selectie- en testaanpak: methodologische literatuur rond exploitselectie en -integratie in gecontroleerde omgevingen illustreert dat exploits en scenario’s best systematisch gekozen, gelogd en geëvalueerd worden om reproduceerbaarheid en interpretatie te verbeteren \autocite{FredrikssonLindkvist2024}.

Op basis van deze vaststellingen kan worden geconcludeerd dat er nood is aan praktijkgericht onderzoek dat:
\begin{enumerate}
    \item de beveiligingsstatus van een reële, heterogene IP-camerainfrastructuur analyseert (inclusief configuratie, firmwarestatus en afhankelijkheden);
    \item bekende kwetsbaarheden en configuratierisico’s empirisch valideert via gecontroleerde proof-of-concepttests binnen een realistische netwerkomgeving, met reproduceerbare logging en scenarioselectie;
    \item netwerkintegratie en aanvalsoppervlak evalueert binnen een concrete retailcontext, rekening houdend met heterogeniteit en beheerrealiteit;
    \item vastgestelde risico’s systematisch vertaalt naar onderbouwde, contextgebonden beveiligingsaanbevelingen door threat modeling (STRIDE) te koppelen aan erkende controlekaders (NIST SP 800-53) en IoT-baselines.
\end{enumerate}

Deze lacune vormt de directe aanleiding voor het voorliggende onderzoek, waarin algemene IoT-beveiligingsinzichten worden toegepast op een specifieke retailomgeving om de praktische impact van camerakwetsbaarheden en netwerkconfiguratie te analyseren en te vertalen naar haalbare mitigaties.